\textbf{COLUNI 2024 -- Questão 35}\par\medskip Leia atentamente o texto a seguir:

\par

\begin{center}\textbf{A Primavera Árabe}\end{center}

\par

No final de 2010, período em que Barack Obama era presidente dos Estados Unidos, ocorreu no Oriente Médio e no norte da África a Primavera Árabe. Esse nome foi atribuído pela mídia ocidental a uma série de protestos ocorridos em países como Tunísia, Egito, Líbia, Argélia, Jordânia, Iêmen, Marrocos e Bahrein. A primeira manifestação ocorreu na Tunísia, quando um comerciante local, após ter sofrido repressão policial, suicidou-se, ateando fogo ao próprio corpo. O fato desencadeou uma série de manifestações em oposição ao governo autoritário do então presidente Zine El Abidine Ben Ali, que se encontrava no poder havia 23 anos. As mobilizações levaram à renúncia do presidente, que fugiu para a Arábia Saudita. Nos meses seguintes, os protestos se expandiram para outros países. Em cada local, havia peculiaridades políticas, mas várias pautas eram comuns, como a situação da população diante da crise econômica, os altos índices de desemprego e a falta de democracia. A principal novidade dessas manifestações foi o uso da internet e das redes sociais para convocar greves, passeatas e comícios.

\par

(KARNAL, Leandro et al. \textbf{Viver História.com}. 9º ano. 1ª edição. São Paulo: Moderna, 2022, p.269)

\par

Com base na leitura do texto e nos conhecimentos sobre a chamada “Primavera Árabe”, assinale a afirmativa CORRETA:\begin{enumerate}
\item O uso das redes sociais virtuais fez com que as pautas e os movimentos se unificassem nas diferentes
        regiões atingidas pela “Primavera Árabe”, permitindo a construção de regimes democráticos.
\item Os resultados da “Primavera Árabe”, até o presente momento, mesmo tendo como ponto de partida
        reivindicações que se referiam à luta contra governos autoritários, não se converteram em realidades
        plenamente democráticas.
\item O fato de todos os países afetados pela “Primavera” serem monarquias islâmicas foi uma das questões
        que possibilitou a expansão dos movimentos contestatórios que eclodiram a partir de 2010, entre o norte
        da África e o Oriente Médio.
\item O Egito, dentre todas as nações atingidas pela “Primavera”, é a única que conseguiu formar um governo
        democrático, após a queda de Hosni Mubarak, que governou o país por 30 anos.\end{enumerate}
